%% Chapter 3: Maximum Flow and Minimum Cut
%% Status: skeleton -- learning goals and section outline fixed, prose to be written.

\chapter{Maximum Flow and Minimum Cut}
\label{ch:max-flow}

Flows are the workhorse of combinatorial optimisation: a large number of
problems that look nothing like a flow problem become easy once they are
modelled as one. We develop the theory, the standard algorithms, and --- most
importantly for the exam --- the modelling tricks.

\section*{Learning goals}
\addcontentsline{toc}{section}{Learning goals}

After this chapter you should be able to:
\begin{itemize}[leftmargin=*]
  \item define flow networks, feasible flows, cuts, and the value of a flow;
  \item prove the max-flow min-cut theorem via residual graphs;
  \item analyse Ford--Fulkerson, Edmonds--Karp and Dinic's algorithm;
  \item reduce flow problems with vertex capacities, several sources or lower bounds to plain maximum flow;
  \item derive Menger's theorem and use it to solve disjoint-path problems.
\end{itemize}

\section{Flow networks}
\label{sec:flow-networks}

\todo{Write this section.}

\section{Residual graphs and augmenting paths}
\label{sec:residual-graphs}

\todo{Write this section.}

\section{The max-flow min-cut theorem}
\label{sec:maxflow-mincut}

\todo{Write this section.}

\section{Algorithms: Ford--Fulkerson, Edmonds--Karp, Dinic}
\label{sec:flow-algorithms}

\todo{Write this section.}
% Planned content: Termination, the integrality theorem, and running-time bounds.

\section{Modelling with flows}
\label{sec:flow-modelling}

\todo{Write this section.}
% Planned content: Vertex capacities, multiple sources and sinks, lower bounds on arcs, and undirected graphs.

\section{Menger's theorem and disjoint paths}
\label{sec:menger}

\todo{Write this section.}

\section{Case study: escaping a grid}
\label{sec:grid-escape}

\todo{Write this section.}
% Planned content: Vertex-disjoint paths from given start points to the boundary, and a careful running-time analysis in terms of the grid size.

\section{Exercises}
\label{sec:ex-max-flow}

\todo{Write this section.}

\begin{poolquestions}
  \poolq{2}{Minimum Vertex Cut}Vertex capacities are handled in \cref{sec:flow-modelling}.
  \poolq{3}{Escaping a Grid}Covered in \cref{sec:grid-escape}.
\end{poolquestions}
